% =============================================================================
%  Phase 4: Modeling & Analysis
%  Rubric: docs/spec/phase4_details.pdf (no page limit).
%  Source of every number: notebooks/modeling.ipynb (executed outputs).
%  Figures: reports/figures/p4_*.png, written by that notebook.
%  Lines marked "% CHECK:" are review notes for the authors, not report text.
%  Build: latexmk -pdf -outdir=build main.tex
% =============================================================================
\documentclass[11pt,letterpaper]{article}
\input{../_template/preamble}
\addbibresource{references.bib}

% ---- Metadata ---------------------------------------------------------------
\projecttitle[Bill Forecasting]{Forecasting Household Month End Electricity Bills from Smart Meter Data}
\phaselabel{Phase 4 \textperiodcentered{} Modeling \& Analysis}
\authors{Akinbobola Akin (3784664), Adejuwon Oluwasegun Isaac (3806766), Edehi Andrew Igue (3778767)}
\course{TME6014 -- Data Analytics}
\reportdate{November 24, 2026}
% CHECK: set to the actual submission date.

\begin{document}
\maketitleblock

% =============================================================================
\section{Introduction and Scope}
\label{sec:intro}

This phase builds and evaluates the models that answer the research question set in Phase~1:

\begin{quote}
\textbf{On day $d$ of a billing month, how accurately can a household's month-end bill be projected, with a
prediction interval, from its smart meter readings? And when should the household be alerted?}
\end{quote}

The task is \textbf{regression}. On day $d$ the kWh used on days $1..d$ is already on the meter, so every model
predicts only the remaining kWh (days $d{+}1$ to month end) and adds it to the known part. The bill is a
\emph{computed proxy}, not a real invoice: kWh $\times$ \pounds 0.14228, the published standard
tariff~\cite{datastore_smartmeters}, with no standing charge. As decided in Phase~3, only standard-tariff (Std)
households are modelled.

The input is the Phase~3 modelling dataset \code{features\_\allowbreak day\_d.parquet}: 254,421 rows (one per
household $\times$ month $\times$ $d$, with $d \in \{7, 14, 21\}$) from 4,428 households, December 2011 to January
2014~\cite{kaggle_smartmeters}. Every number and figure in this report comes from the notebook
\code{notebooks/modeling.ipynb} (Section~\ref{sec:code}).

% =============================================================================
\section{Data Splitting}
\label{sec:split}

The data are a time series, so a random split would let a model learn from months that come after the ones it is
scored on. The split is therefore \textbf{by calendar month}. Every row of a household-month (its three $d$
values) always lands in the same part, so the model never sees the same month at another $d$. All choices
(features, hyperparameters, interval method, alert rule) are made on three expanding-window validation folds. The
last four months are a test set, used once, in Section~\ref{sec:test} (Table~\ref{tab:split}).

\begin{table}[!ht]
  \centering
  \caption{Time-based split. Each fold trains on every month before its validation window. Lag availability
  is the share of rows where the feature exists.}
  \label{tab:split}
  \small
  \begin{tabularx}{\linewidth}{L l r r r r}
    \tableheadrow
    \colhead{Part} & \colhead{Evaluated on} & \colhead{Training rows} & \colhead{Evaluated rows} &
    \colhead{Has last month} & \colhead{Has same month LY} \\
    \tableheadend
    Fold 1 & Jan--Mar 2013 & 92,247 & 38,676 & 99.6\% & 11.6\% \\
    Fold 2 & Apr--Jun 2013 & 130,923 & 37,983 & 99.7\% & 42.0\% \\
    Fold 3 & Jul--Sep 2013 & 168,906 & 37,194 & 99.7\% & 81.7\% \\
    Test & Oct 2013--Jan 2014 & 206,100 & 48,321 & 99.6\% & 96.7\% \\
    \bottomrule
  \end{tabularx}
\end{table}

Most households joined in mid-2012, so ``same month last year'' exists for only 0.2\% of the rows that train
fold~1, but for 96.7\% of test rows. Preprocessing (imputation, scaling) sits inside a scikit-learn
pipeline~\cite{sklearn}, so it is fitted on each fold's training months only. The gradient boosting model's
built-in early stopping is switched off, because it holds out a \emph{random} 10\% of rows, which would break the
time split.

% =============================================================================
\section{Models Tested}
\label{sec:models}

\subsection{Benchmarks}
The three benchmarks promised in Phase~1 need no training:
\begin{itemize}
  \item \textbf{Run-rate}: to-date daily mean $\times$ days in the month.
  \item \textbf{Last month}: the previous month's full bill.
  \item \textbf{Same month last year}: that month's bill one year earlier.
\end{itemize}
A fallback \textbf{chain} (same month last year, else last month, else run-rate) gives every row a forecast.

\subsection{Candidate models}
\begin{itemize}
  \item \textbf{Ridge regression}, the interpretable linear reference. Missing lags are imputed with the median,
  with indicator columns, and features are standardised. \emph{Strength:} simple, fast, and its coefficients can
  be read. \emph{Weakness:} it cannot represent interactions such as ``early in the month and a large household''
  without hand-made terms, and it fits the mean, so a few very large homes pull it around.
  \item \textbf{Histogram gradient boosting} (scikit-learn's \code{HistGradientBoostingRegressor}, in the style of
  LightGBM~\cite{ke2017lightgbm}). \emph{Strength:} it learns non-linear effects and interactions, handles missing
  lags natively, trains in seconds on 250,000 rows, and supports absolute-error and quantile losses.
  \emph{Weakness:} less transparent (we use permutation importance to explain it), and it can overfit calendar
  features when history is short (Section~\ref{sec:features}).
  \item \textbf{Per-household model}, for the secondary question (pooled versus per-household). A regression per
  household is not realistic: a household has a median of about 20 months, so about 20 rows per $d$. Instead, each
  household's own average correction at the same $d$ over its \emph{earlier} months is shrunk toward the run-rate
  with weight $k/(n+k)$, where $n$ is the number of earlier months. $k = 0$ trusts the household's history fully,
  and a large $k$ falls back to the run-rate.
\end{itemize}

\subsection{Target transformation}
\label{sec:target}
Errors grow with household size: a 5\% miss is about \pounds 1 for a small home and \pounds 10 for a large one.
Besides the raw target (\code{remaining\_kwh}), we tested a \textbf{ratio target}:
\[
  \text{ratio} = \frac{\text{remaining kWh}}{\text{scale}}, \qquad
  \text{scale} = \max(\text{to-date daily mean},\ 0.5) \times \text{days remaining}.
\]
A ratio of 1 means the run-rate is exactly right, so the model only has to learn the correction to it. The 0.5
kWh/day floor (415 rows) stops near-empty homes from producing huge ratios. The model is trained with an
absolute-error loss weighted by the scale, which makes the loss exactly the absolute error in kWh, and so in
pounds. Four relative features follow the same idea:
\begin{itemize}
  \item the last 7 days against the to-date mean;
  \item last month against the floored run-rate;
  \item the 3-month mean against the floored run-rate;
  \item same month last year against the floored run-rate.
\end{itemize}

\subsection{Prediction intervals}
Three split-conformal methods~\cite{lei2018conformal} build a 90\% interval around the gradient boosting forecast,
each calibrated separately for each $d$:
\begin{itemize}
  \item \textbf{absolute}: forecast $\pm$ a fixed amount in pounds;
  \item \textbf{scaled}: forecast $\pm$ a multiple of the scale, so the width grows with usage and days left;
  \item \textbf{CQR} (conformalized quantile regression~\cite{romano2019cqr}): two extra gradient boosting models
  predict the 5th and 95th percentiles, then a conformal correction adjusts them.
\end{itemize}
Calibration uses out-of-fold errors, so each error comes from a model that never saw that month. No interval
falls below the bill already on the meter.

% =============================================================================
\section{Evaluation Metrics}
\label{sec:metrics}

All point metrics are computed on the month-end bill in pounds:
\begin{itemize}
  \item \textbf{MAE}, the main metric: the household's typical error in pounds;
  \item \textbf{RMSE}, which penalises large misses;
  \item \textbf{MdAPE}, the median absolute percentage error, which is scale-free;
  \item \textbf{bias}, the mean of forecast minus actual, where negative means under-forecasting.
\end{itemize}
Intervals are judged by \textbf{coverage} (target 90\%) and mean \textbf{width}. Coverage is checked overall and
within groups: by $d$, by month, and by household size (terciles of the run-rate bill). This matters because an
interval can hit 90\% on average while failing large households. Alerts are scored as a classification:
\textbf{precision} (share of alerts that were right), \textbf{recall} (share of real overruns that got an alert)
and \textbf{F1}. Whether the gap to the run-rate is real is tested with a household-block bootstrap (1,000
resamples of whole households), which keeps each household's months and $d$ values together.

% =============================================================================
\section{Model Comparison and Optimization}
\label{sec:validation}

\subsection{Benchmarks and default models}
Table~\ref{tab:cand} shows mean validation results across the three folds. The run-rate (\pounds 3.43) is far
ahead of the history-only benchmarks. Last month (\pounds 6.83), same month last year (\pounds 8.85, available on
45\% of rows) and the chain (\pounds 7.27) all predict the whole month and ignore the $d$ days already on the
meter. On the same rows the run-rate scores about \pounds 3.43, so their losses are not caused by which rows they
cover.

With default settings, none of the four candidates beat the run-rate. For gradient boosting:
\begin{itemize}
  \item the absolute-error loss beats the squared loss (\pounds 3.49 vs 3.61), because the squared loss chases a
  few very large months;
  \item the ratio target matches the raw target on MAE and has the lowest RMSE (7.51), so it was kept.
\end{itemize}

\begin{table}[!ht]
  \centering
  \caption{Mean validation results across the three folds (\pounds), default hyperparameters.}
  \label{tab:cand}
  \small
  \begin{tabularx}{\linewidth}{L r r r r}
    \tableheadrow
    \colhead{Model} & \colhead{MAE} & \colhead{RMSE} & \colhead{MdAPE \%} & \colhead{Bias} \\
    \tableheadend
    Run-rate benchmark & 3.43 & 7.72 & 4.08 & $-0.18$ \\
    Last month benchmark & 6.83 & 13.29 & 10.13 & 1.97 \\
    Same month last year benchmark (45\% of rows) & 8.85 & 17.48 & 11.87 & $-1.30$ \\
    Chain benchmark & 7.27 & 13.91 & 11.04 & 1.57 \\
    \midrule
    Ridge, raw kWh target & 3.81 & 7.69 & 6.26 & 0.17 \\
    Gradient boosting, raw target, squared loss & 3.61 & 7.89 & 4.65 & $-0.44$ \\
    Gradient boosting, raw target, absolute loss & 3.49 & 7.81 & 4.42 & $-0.36$ \\
    Gradient boosting, ratio target, absolute loss & 3.49 & 7.51 & 4.43 & $-0.62$ \\
    \bottomrule
  \end{tabularx}
\end{table}

\subsection{Feature selection}
\label{sec:features}
Features were tested in related groups, because related features stand in for each other, so dropping one at a
time understates their value. A group is dropped if removing it lowers the mean validation MAE by more than
\pounds 0.05 (Table~\ref{tab:ablation}).

\begin{table}[!ht]
  \centering
  \caption{Feature-group ablation, gradient boosting with the ratio target. Change is mean MAE versus the full
  feature set (negative means dropping the group helps).}
  \label{tab:ablation}
  \small
  \begin{tabularx}{\linewidth}{L r r r r r}
    \tableheadrow
    \colhead{Group dropped} & \colhead{Fold 1} & \colhead{Fold 2} & \colhead{Fold 3} & \colhead{Mean} &
    \colhead{Change} \\
    \tableheadend
    None (all features) & 4.473 & 3.412 & 2.463 & 3.449 & 0.000 \\
    Month of year (sin/cos) & 3.849 & 3.576 & 2.511 & 3.312 & $-0.137$ \\
    Temperature to date & 4.511 & 3.393 & 2.446 & 3.450 & 0.001 \\
    Earlier-month lags & 4.377 & 3.368 & 2.465 & 3.404 & $-0.046$ \\
    Relative features & 4.524 & 3.424 & 2.514 & 3.487 & 0.038 \\
    ACORN group & 4.478 & 3.394 & 2.462 & 3.445 & $-0.005$ \\
    Rest-of-month calendar & 4.498 & 3.790 & 2.464 & 3.584 & 0.135 \\
    \bottomrule
  \end{tabularx}
\end{table}

Only the month-of-year pair was dropped. Removing it cuts fold~1 MAE from 4.47 to 3.85. With one or two years of
history, the model sees each calendar month only once or twice, so it learns that year's weather, not a seasonal
pattern. Nineteen features remain.

Permutation importance on fold~3 (20,000 sampled rows) confirms where the signal is: the relative features.
Shuffling a feature raises the bill MAE by:
\begin{itemize}
  \item last month against the run-rate: \pounds 0.24;
  \item the last 7 days against the to-date mean: \pounds 0.15;
  \item same month last year against the run-rate: \pounds 0.07;
  \item the 3-month mean against the run-rate: \pounds 0.05;
  \item each raw size feature: at most \pounds 0.013.
\end{itemize}
The raw size features matter little once the ratio target has removed household size.
% CHECK: temperature and weekend days remaining have slightly negative importance (-0.009); kept because the
% group test showed changes under the 0.05 rule.

\subsection{Hyperparameter tuning}
\paragraph{Ridge} The penalty $\alpha \in \{0.1, 1, 10, 100, 1000\}$ was tuned. The best is $\alpha = 10$ (mean
MAE \pounds 3.675). $\alpha = 100$ is nearly identical, and $\alpha = 1000$ is worse (\pounds 3.797).

\paragraph{Gradient boosting} A grid of eight configurations was tested:
\begin{itemize}
  \item learning rate $\{0.03, 0.1\}$;
  \item leaves per tree $\{7, 31\}$;
  \item minimum rows per leaf $\{200, 2000\}$;
  \item 300 iterations in every case.
\end{itemize}
The best is learning rate 0.03, 31 leaves and 200 rows per leaf (mean MAE \pounds 3.271, RMSE 7.30). The top five
configurations are within \pounds 0.03 of each other, so the result does not hinge on the exact setting. Tuning
and feature selection together cut gradient boosting's mean MAE from \pounds 3.49 to \pounds 3.27.

\subsection{Tuned models by fold}
After tuning, gradient boosting beats the run-rate in \emph{every} fold. Ridge loses to the run-rate in every fold,
so it is kept only as a reference (Table~\ref{tab:folds}).

\begin{table}[!ht]
  \centering
  \caption{Validation MAE (\pounds) of the tuned models by fold.}
  \label{tab:folds}
  \small
  \begin{tabularx}{0.7\linewidth}{L r r r r}
    \tableheadrow
    \colhead{Model} & \colhead{Fold 1} & \colhead{Fold 2} & \colhead{Fold 3} & \colhead{Mean} \\
    \tableheadend
    Run-rate & 3.928 & 3.646 & 2.702 & 3.425 \\
    Ridge ($\alpha = 10$) & 4.259 & 4.055 & 2.711 & 3.675 \\
    Gradient boosting (tuned) & \textbf{3.783} & \textbf{3.544} & \textbf{2.488} & \textbf{3.271} \\
    \bottomrule
  \end{tabularx}
\end{table}

\subsection{Pooled versus per-household}
The per-household model gets \emph{better} the less it trusts the household (Figure~\ref{fig:pooled}):
\begin{itemize}
  \item at $k = 0$ (full trust) the mean MAE is \pounds 3.92;
  \item at $k = 128$ it is \pounds 3.43, which is simply the run-rate.
\end{itemize}
A household's past ``rest of month versus first $d$ days'' correction does not carry over to the next month.
Adding this household history to the pooled model as features also makes it slightly worse (\pounds 3.32 vs 3.27).
\textbf{The pooled model wins.} The pattern it learns (households running above their usual level tend to keep
running above it) is shared across households, and 20 months per household is too few to learn it separately.

\begin{figure}[!ht]
  \centering
  \includegraphics[width=0.75\linewidth]{p4_pooled_vs_household.png}
  \caption{Mean validation MAE: per-household models with shrinkage $k$, pooled gradient boosting with and
  without household history features, and the run-rate.}
  \label{fig:pooled}
\end{figure}

\subsection{Choosing the interval method}
The method was chosen on fold~3 using intervals calibrated on folds~1--2. The rule: pick the narrowest method that
covers at least 87\% in every group (Table~\ref{tab:ivfold3}).
\begin{itemize}
  \item \textbf{Absolute} over-covers overall (94.8\%) but falls to 88.7\% for large households, because a fixed
  width is too wide for small bills and too narrow for large ones.
  \item \textbf{Scaled} fixes large households but under-covers small ones (87.4\%).
  \item \textbf{CQR} stays between 88\% and 92\% in every group and is the narrowest, so it was chosen.
\end{itemize}

\begin{table}[!ht]
  \centering
  \caption{90\% interval methods on fold 3 (calibrated on folds 1--2).}
  \label{tab:ivfold3}
  \small
  \begin{tabularx}{0.8\linewidth}{L r r r}
    \tableheadrow
    \colhead{Method} & \colhead{Coverage} & \colhead{Worst group} & \colhead{Mean width (\pounds)} \\
    \tableheadend
    Absolute & 94.8\% & 88.7\% & 16.55 \\
    Scaled & 91.0\% & 87.4\% & 12.72 \\
    CQR & 90.5\% & 88.3\% & \textbf{10.54} \\
    \bottomrule
  \end{tabularx}
\end{table}

% =============================================================================
\section{Final Test Results}
\label{sec:test}

The tuned models were refitted on all months up to September 2013 and scored once on October 2013 to January 2014
(48,321 rows from 4,118 households).

\subsection{Point forecasts}
Gradient boosting has the lowest error (Table~\ref{tab:test}):
\begin{itemize}
  \item MAE \pounds 3.60 against \pounds 3.80 for the run-rate, about 5\% lower;
  \item MdAPE 4.2\%;
  \item the household-block bootstrap puts the MAE difference at $-$\pounds 0.198 (95\% CI $-0.234$ to $-0.157$),
  so the gain is real.
\end{itemize}
By forecast day (Table~\ref{tab:testd} and Figure~\ref{fig:testerr}), the gain is largest at $d = 7$, where the
run-rate has least to go on, and clear at $d = 21$. At $d = 14$ the model and the run-rate are
\textbf{indistinguishable}. Ridge is no better than the run-rate overall.

\begin{table}[!ht]
  \centering
  \caption{Test results, October 2013 to January 2014 (\pounds).}
  \label{tab:test}
  \small
  \begin{tabularx}{\linewidth}{L r r r r r}
    \tableheadrow
    \colhead{Model} & \colhead{Available} & \colhead{MAE} & \colhead{RMSE} & \colhead{MdAPE \%} & \colhead{Bias} \\
    \tableheadend
    Run-rate & 100\% & 3.80 & 8.60 & 4.21 & $-0.78$ \\
    Last month & 99.6\% & 7.21 & 14.84 & 9.29 & $-3.56$ \\
    Same month last year & 96.7\% & 9.84 & 21.53 & 12.35 & 2.33 \\
    Chain & 100\% & 9.72 & 21.29 & 12.25 & 2.12 \\
    Ridge & 100\% & 3.83 & 8.18 & 4.97 & $-1.43$ \\
    Gradient boosting & 100\% & \textbf{3.60} & \textbf{8.06} & 4.24 & $-1.28$ \\
    \bottomrule
  \end{tabularx}
\end{table}

\begin{table}[!ht]
  \centering
  \caption{Test MAE by forecast day (\pounds) and the bootstrap 95\% CI for gradient boosting minus run-rate.}
  \label{tab:testd}
  \small
  \begin{tabularx}{\linewidth}{L r r r Y{1}}
    \tableheadrow
    \colhead{$d$} & \colhead{Run-rate} & \colhead{Ridge} & \colhead{Gradient boosting} &
    \colhead{Difference (95\% CI)} \\
    \tableheadend
    7 & 5.534 & 5.329 & \textbf{5.083} & $-0.451$ ($-0.516$ to $-0.379$) \\
    14 & 3.609 & 3.633 & \textbf{3.594} & $-0.015$ ($-0.053$ to $+0.027$) \\
    21 & 2.248 & 2.516 & \textbf{2.120} & $-0.128$ ($-0.155$ to $-0.101$) \\
    \bottomrule
  \end{tabularx}
\end{table}

\begin{figure}[!ht]
  \centering
  \includegraphics[width=\linewidth]{p4_test_error.png}
  \caption{Test MAE by forecast day (left) and mean bias by test month (right) for the run-rate, Ridge and gradient
  boosting.}
  \label{fig:testerr}
\end{figure}

\paragraph{Bias} Every model under-forecasts in October and November 2013. Gradient boosting's bias is
$-$\pounds 1.63 and $-$\pounds 2.78 in those months, because usage rises through the month as days get shorter and
colder. Gradient boosting is more biased overall than the run-rate ($-$\pounds 1.28 vs $-$\pounds 0.78). Its
absolute-error loss targets the \emph{median} remaining usage, which is below the mean when a few households
spike. This trade-off gives a lower typical error at the cost of a small systematic under-forecast.

\subsection{Prediction intervals}
Calibrated on the out-of-fold errors of all three folds (January to September 2013, including winter months), the
CQR intervals cover \textbf{88.9\%} of test bills (target 90\%), with a mean width of \pounds 14.44
(Figure~\ref{fig:cov}). The absolute method is wider (\pounds 14.83) and covers less (88.1\%). It covers only
74.3\% of large households, where CQR covers 87.1\%. CQR's weakest groups are November 2013 (86.3\%) and $d = 21$
(86.6\%), the cases where every model under-forecasts.

\begin{figure}[!ht]
  \centering
  \includegraphics[width=\linewidth]{p4_interval_coverage.png}
  \caption{Test coverage of the three 90\% interval methods by forecast day, month and household size (left),
  and mean width (right).}
  \label{fig:cov}
\end{figure}

% =============================================================================
\section{Alert Decision Layer}
\label{sec:alerts}

\textbf{Rule:} on day $d$, alert the household if the projected bill exceeds its baseline by more than $X$.
\textbf{Truth:} the actual bill exceeded the baseline by more than $X$. Three choices were compared, each made on
fold~3 and reported on the test months.

\paragraph{Baseline}
\begin{itemize}
  \item \textbf{Trailing 3-month mean}: labels 35.2\% of test bills as overruns, against 14--15\% for the other two
  baselines. In autumn and winter it compares each month with cheaper months, so it flags the normal seasonal rise.
  \item \textbf{Same month last year}: the fairest across seasons, but missing for 18\% of fold~3 rows (households
  less than a year old).
  \item \textbf{Chosen:} same month last year where it exists, otherwise the trailing 3-month mean. It is
  available for 99.8\% of rows.
\end{itemize}

\paragraph{Forecast day} Alert quality rises with $d$: test F1 for the model rule at $X = 20\%$ is 0.67, 0.77 and
0.86 for $d = 7$, 14 and 21. But at $d = 21$ only about a week is left to change behaviour, and at $d = 7$ the
model rule misses 43\% of overruns. \textbf{$d = 14$ is chosen}: half the month is left to act, with F1 near 0.8.

\paragraph{Threshold $X$} $X$ is the household's tolerance, not something to optimise, and alert quality barely
changes between 10\% and 30\% (Figure~\ref{fig:alerts}). The default is \textbf{20\%}.

\paragraph{Alert rule} Three triggers were compared (Table~\ref{tab:alerts}):
\begin{itemize}
  \item the run-rate bill;
  \item the gradient boosting forecast;
  \item the lower bound of the 90\% interval.
\end{itemize}
On fold~3 the run-rate rule has a slightly higher F1 than the model rule (0.785 vs 0.778), so it is chosen. On the
test months it flags 13.4\% of household-months, which equals the true overrun rate, with 80\% precision and 80\%
recall. The model rule trades recall for precision (85.5\% / 70.6\%), because its downward bias makes it alert less
often. The lower-bound rule is almost never wrong (98.4\% precision) but catches less than a third of overruns, so
it suits a ``near-certain overrun'' message rather than the main alert.

\begin{table}[!ht]
  \centering
  \caption{Alert rules at $d = 14$, $X = 20\%$, baseline: same month last year, else trailing 3-month mean.}
  \label{tab:alerts}
  \small
  \begin{tabularx}{\linewidth}{L r r r r r}
    \tableheadrow
    \colhead{Rule} & \colhead{Fold 3 F1} & \colhead{Test alert \%} & \colhead{Test precision} &
    \colhead{Test recall} & \colhead{Test F1} \\
    \tableheadend
    Run-rate (chosen) & \textbf{0.785} & 13.4 & 0.797 & \textbf{0.797} & \textbf{0.797} \\
    Model point forecast & 0.778 & 11.1 & 0.855 & 0.706 & 0.774 \\
    Model interval lower bound & 0.415 & 4.3 & \textbf{0.984} & 0.313 & 0.475 \\
    \bottomrule
  \end{tabularx}
\end{table}

\begin{figure}[!ht]
  \centering
  \includegraphics[width=\linewidth]{p4_alerts.png}
  \caption{Test precision and recall of the three alert rules against the threshold $X$, at $d = 14$.}
  \label{fig:alerts}
\end{figure}

% =============================================================================
\section{Final Approach and Justification}
\label{sec:final}

\begin{enumerate}
  \item \textbf{Forecast:} pooled histogram gradient boosting on the ratio target, trained with an absolute-error
  loss.
  \begin{itemize}
    \item \emph{Settings:} learning rate 0.03, 31 leaves, at least 200 rows per leaf, 300 iterations, 19 features.
    \item \emph{Why:} it was the only model to beat the run-rate in every validation fold, and it does so on the
    test months, by a statistically clear margin.
    \item \emph{What made the gain:} predicting the correction to the run-rate, plus features comparing the month
    with the household's own recent level. Model complexity contributed little.
  \end{itemize}
  \item \textbf{Uncertainty:} CQR 90\% intervals, chosen because they keep coverage near 90\% across forecast days,
  months and household sizes while being the narrowest.
  \item \textbf{Alert:} $d = 14$, $X = 20\%$ by default, baseline same month last year else trailing 3-month mean,
  triggered by the run-rate threshold. The model's forecast and interval are shown inside the alert message (for
  example, ``projected \pounds 62, likely between \pounds 55 and \pounds 70''). For the yes/no decision a simple
  threshold is as good as the model. The model's value is a more accurate number with honest uncertainty.
\end{enumerate}

The claim is deliberately modest. The run-rate is a strong benchmark, and gradient boosting improves on it by
about 5\% overall. At the chosen $d = 14$ the point forecast is no better than the run-rate. At that day the
model's contribution is the calibrated interval, not a lower error.
% CHECK: the team may prefer d = 21 if accuracy matters more than lead time; all numbers for d = 21 are in the
% notebook (section 9).

% =============================================================================
\section{Limitations}
\label{sec:limits}

\begin{itemize}
  \item \textbf{Bill proxy.} kWh $\times$ 14.228~p/kWh, with no standing charge. This is not a real invoice.
  \item \textbf{Std households only.} Time-of-use households remain the stretch goal.
  \item \textbf{Complete months only.} Phase~2 kept only complete household-months, but a live system would not
  know in advance that a month will be complete.
  \item \textbf{Hardest-season test window.} The test months are four autumn and winter months, when usage rises
  and every model under-forecasts. Errors in spring and summer are expected to be lower, as folds~2 and~3 suggest.
  \item \textbf{About two years of history.} This is too little to learn seasonality reliably, and same-month-last-
  year lags exist only from 2013.
  \item \textbf{Conformal guarantee.} It assumes past errors resemble future ones (exchangeability), which seasonal
  data does not strictly satisfy~\cite{hyndman2021fpp3}. Coverage was therefore checked empirically, and an
  unusual month such as a cold snap can still break it.
\end{itemize}

% =============================================================================
\section{Source Code and Reproducibility}
\label{sec:code}

All modelling code is in \code{notebooks/modeling.ipynb}, which passes Restart \& Run All in about ten minutes:
\begin{itemize}
  \item \emph{Input:} it reads \code{data/processed/features\_\allowbreak day\_d.parquet} with DuckDB.
  \item \emph{Configuration:} every path, split date, grid and threshold is set in its first cell.
  \item \emph{Random seed:} fixed (42).
  \item \emph{Outputs:} the test forecasts, intervals and alerts (\code{predictions\_test.parquet} in
  \code{data/processed/}) and the figures (\code{p4\_*.png} in \code{reports/figures/}).
\end{itemize}
A final cell runs nine checks and stops the notebook if any fails:
\begin{itemize}
  \item no training or calibration row from the test months;
  \item fold training windows end before their validation windows;
  \item every test row has a forecast and an interval;
  \item no lower bound lies above its upper bound;
  \item no forecast falls below the bill already on the meter;
  \item the saved rows match the test set.
\end{itemize}
% CHECK: add the repository link if the course asks for it.

\printbibliography[title={References}]

\end{document}
